\chapter[Constructibility and Forcing]{Constructibility, Forcing, and Relative Consistency}
\label{ch:constructibility-forcing}

\section{The constructible hierarchy}
\label{sec:constructibility-forcing-01}

\begin{definition}[Definable subsets]
\label{def:p06-def-empty-base}
Set
\[
  \operatorname{Def}(\varnothing)\defeq\{\varnothing\}.
\]
For a nonempty set \(A\), let \(\operatorname{Def}(A)\) be the collection of subsets of
\(A\) definable over the structure \((A,\in)\) by a first-order formula with
parameters from \(A\):
\[
  \operatorname{Def}(A)
  \defeq
  \set{
    X\subseteq A:
    \begin{gathered}
      \text{for some formula }\varphi(y,\bar z)
      \text{ and }\bar a\in A^{<\omega},\\
      X=\set{x\in A:(A,\in),s_{x,\bar a}\entails\varphi}
    \end{gathered}
  }.
\]
Here \(A^{<\omega}\) denotes the set of finite tuples from \(A\), and
\(s_{x,\bar a}\) maps \(y\) to \(x\) and each displayed variable in
\(\bar z\) to the corresponding entry of \(\bar a\); its values on other
variables are irrelevant.
The satisfaction clause applies to nonempty \(A\).  The separate empty-base
clause makes the unique subset of \(\varnothing\) the sole member of
\(\operatorname{Def}(\varnothing)\), without evaluating a formula in an
empty structure.
The satisfaction relation in this definition is the set-sized satisfaction
relation constructed in the metatheory.  It is not a truth predicate internal
to \(\ZF\) for the entire universe.
\end{definition}

\begin{definition}[Constructible hierarchy]
Define by transfinite recursion
\[
  L_0=\varnothing,\qquad
  L_{\alpha+1}=\operatorname{Def}(L_\alpha),\qquad
  L_\lambda=\bigcup_{\beta<\lambda}L_\beta
\]
for limit \(\lambda\), and write
\[
  L=\bigcup_{\alpha\in\operatorname{Ord}}L_\alpha.
\]
A set is constructible iff it belongs to \(L\).
In particular, \(L_1=\operatorname{Def}(\varnothing)=\{\varnothing\}\).
\end{definition}

The ordinary Tarskian satisfaction relation fixed earlier takes a set-sized
structure on its left and therefore cannot take the proper class \(L\) as an
argument.  To state the next result without adding class structures or a class
truth predicate, use the following formula-by-formula abbreviation.  For a
standard formula \(\varphi\) of \(\Lang_{\in}\), let
\(\varphi^L\) be its metalinguistic relativization to the displayed definable
class: atomic formulas are unchanged, the operation commutes with the
propositional connectives, and
\[
  (\exists x\,\psi)^L\defeq
    \exists x\,(x\in L\land\psi^L),
  \qquad
  (\forall x\,\psi)^L\defeq
    \forall x\,(x\in L\to\psi^L).
\]
Here \(x\in L\) is eliminated using the formula defining constructibility.
Accordingly, \(L\entails\sigma\) abbreviates the ambient assertion
\(\sigma^L\), and \(L\entails T\) means that this holds for every standard
sentence, including every displayed schema instance, represented by the
theory expression \(T\).  This is a formula-by-formula metalinguistic
convention.  It neither treats the proper class \(L\) as a set-sized structure
nor posits a single truth predicate for all formula codes.

\begin{theorem}[The constructible class L]
Assume \(\ZF\) in the ambient universe.  The class \(L\) is transitive and
contains every ordinal.
\end{theorem}

\paragraph{Complete proof.}
The full derivation is supplied in
\cref{sec:p06-constructible-hierarchy-foundations} and
\cref{annex:p06.constructible-class-foundation}. It establishes the
ordinal, recursion and satisfaction prerequisites, proves transitivity
of every \(L_\alpha\), and gives the defining formula that places
\(\alpha\) in \(L_{\alpha+1}\); see
\cref{thm:p06-constructible-level-height}. See also \cite{Jech2003}.

\paragraph{Cited theorem (relativization scheme for the constructible
class).}
For every axiom \(\sigma\) of \(\ZFC+\mathsf{GCH}+(V=L)\), including every
instance of the separation and replacement schemas,
\[
  \ZF\derives\sigma^{L},
\]
where \((\cdot)^{L}\) is the relativization convention displayed above.
Each instance of the scheme is a theorem of \(\ZF\) about the definable
class \(L\); the proofs---reflection to set-sized levels, the canonical
definable well-order of \(L\), and condensation---are supplied by
\cite[Chapter~13, pp.~175--200]{Jech2003}.  No satisfaction relation and no
consistency implication is asserted here; the exact relative-consistency
statements are derived from this scheme in the final section of this
chapter.
